\documentclass[a4paper, 12pt]{article}
% math symbols
\usepackage{amssymb}
\usepackage{amsmath}
\usepackage{mathrsfs}
\usepackage{mathseries}


\usepackage[margin = 2cm]{geometry}

\tolerance = 1000
\emergencystretch = 0.74cm



\pagestyle{empty}
\parindent = 0mm

\renewcommand{\coursetitle}{DM/ML}
\setcounter{curtask}{1}

\setmathstyle{\includegraphics[scale = 0.05]{../pics/utia-rest.png}}{Теория информации}{2 курс}


\begin{document}

\libproblem*{cc}{chrom-number}
\libproblem{cc}{KW-hard-exists}
\dzcomment{
    \textit{Подсказка:} формульная сложность такой функции будет $2^{\Omega(n)}$.
}

\begin{definition*}
    \hinttext{Неформально.} Вероятностным коммуникационным протоколом будем называть протокол $\Pi$, в
    котором у Алисы и Боба есть доступ к общим (т.е. оба игрока видят случайные биты) случайным битам. Их
    цель найти значение функции $f(x, y)$ при этом:
    $$
        \forall x, y ~~ \Pr\limits_{r}[\Pi(x, y) \neq f(x, y)] \le \varepsilon,
    $$
    для некоторого параметра $\varepsilon$.

    Минимальное число бит, которым нужно обменяться Алисе и Бобу для того, чтобы посчитать значение
    функции с указанными ограничениями будем обозначать $\Rpub[\varepsilon]$.
\end{definition*}

\libproblem{cc}{rand-eq}
\libproblem{cc}{balance}
\libproblem{cc}{clique-ind}

\begin{definition*}
    Пусть $f\colon X \times Y \to Z$ и $\mu$~--- распределение на $X \times Y$. Заметим, что для любого
    коммуникационного протокола $\Pi$ для функции $f$ распределение $\mu$ индуцирует распределение на
    листьях данного протокола естественным образом. \deftext{Внешней информационной стоимостью} (или
    \deftext{внешним информационным разглашением}) протокола $\Pi$ по распределению $\mu$ будем называть
    величину:
    $$\ICext[\mu](\Pi) \coloneqq I(\Pi(X, Y) : X, Y).$$
    Также определим внешнюю информационную сложность самой функции
    $\ICext[\mu](f) \coloneqq \min\limits_{\Pi} \ICext[\mu](\Pi)$.

    \deftext{Внутренней информационной стоимостью} протокола $\Pi$ по распределению $\mu$ будем называть
    величину:
    $$\ICext[\mu](\Pi) \coloneqq I(\Pi(X, Y) : X \mid Y) + I(\Pi(X, Y) : Y \mid X).$$
\end{definition*}

\libproblem{cc}{inner-inf-upper}
\libproblem{cc}{rand-gt}

\dzcomment{Пункт б сложный, пункт а будет оцениваться отдельно.}

\begin{definition*}
    \deftext{Идеальная схема разделения секрета}~--- это совершенная схема разделения секрета с
    дополнительным требованием <<экономности>>.
    $$
        \forall i \in \{1, 2, \dots, n\},\ h(S_i) \le h(S_0).
    $$
\end{definition*}

\libproblem{inf-theory}{secret-not-ideal}

\end{document}



%%% Local Variables:
%%% mode: latex
%%% TeX-master: t
%%% End:
